\subsection{Threat Modeling dengan STRIDE}

Analisis ancaman pada penelitian ini dilakukan menggunakan metode STRIDE untuk mengidentifikasi ancaman keamanan berdasarkan elemen-elemen pada Data Flow Diagram (DFD). Penerapan STRIDE pada penelitian ini menggunakan pendekatan STRIDE-per-Element, yaitu pendekatan yang memetakan setiap kategori ancaman terhadap jenis elemen DFD yang relevan. Pemetaan tersebut mengacu pada klasifikasi STRIDE-per-Element Fig.~\ref{}, sehingga setiap ancaman dianalisis sesuai dengan jenis elemen yang mengalaminya.

Berdasarkan pemetaan STRIDE-per-Element tersebut, penelitian ini menganalisis lima kategori ancaman, yaitu Spoofing, Tampering, Repudiation, Information Disclosure, dan Elevation of Privilege. Sementara itu, kategori Denial of Service (DoS) tidak dianalisis karena penelitian ini berfokus pada perancangan kebutuhan audit trail untuk mendukung keamanan dan investigasi insiden, bukan pada aspek performa, ketersediaan layanan (availability), atau ketahanan sistem terhadap serangan yang bertujuan mengganggu operasional layanan.

Hasil analisis STRIDE digunakan untuk mengidentifikasi potensi ancaman pada setiap elemen DFD sehingga dapat menjadi dasar dalam menentukan kebutuhan audit trail yang diperlukan untuk mendukung proses deteksi, investigasi, dan pembuktian insiden keamanan pada sistem DecMed. Sebagai penanda unik antar-ancaman dan untuk mempermudah penelusuran keterkaitan antara ancaman STRIDE dengan kebutuhan event audit pada subbab berikutnya, setiap baris ancaman diberikan Kode Ancaman dengan format \texttt{[Kategori]-[Nomor Urut]}, misalnya \texttt{SP-01} untuk ancaman \textit{Spoofing} pertama, \texttt{TP-01} untuk ancaman \textit{Tampering} pertama, dan seterusnya.



% % STRIDE Spoofing
\subsubsection{Spoofing}

Berdasarkan pemetaan STRIDE-per-Element, kategori Spoofing dianalisis pada elemen External Entity dan Process. Analisis ini bertujuan untuk mengidentifikasi kemungkinan pihak yang tidak berwenang memalsukan identitas entitas atau layanan sehingga dapat memperoleh akses atau menjalankan operasi atas nama entitas yang sah.

\begin{xltabular}{\textwidth}{|>{\centering\arraybackslash}p{1.8cm}|>{\raggedright\arraybackslash}p{3.2cm}|X|}
    \caption{Analisis STRIDE Kategori Spoofing}
    \label{tab:stride-spoofing} \\
    \hline
    \textbf{Kode Ancaman} &
    \textbf{Komponen} &
    \textbf{Deskripsi Ancaman} \\
    \hline
    \endfirsthead

    \multicolumn{3}{c}%
    {} \\
    \hline
    \textbf{Kode Ancaman} &
    \textbf{Komponen} &
    \textbf{Deskripsi Ancaman} \\
    \hline
    \endhead

    \hline
    \multicolumn{3}{r}{} \\
    \endfoot

    \hline
    \endlastfoot

    SP-01 &
    E1 -- Patient &
    Pihak tidak sah mengklaim sebagai pasien tertentu menggunakan kredensial sah yang tercuri . \\
    \hline

    SP-02 &
    E2 -- Hospital Personnel &
    Pihak tidak sah mengklaim sebagai tenaga medis menggunakan kredensial sah yang tercuri. \\
    \hline

    SP-03 &
    E3 -- Ministry Administrator &
    Pihak tidak sah mengklaim sebagai administrator Kementerian Kesehatan menggunakan kredensial sah yang tercuri. \\
    \hline

    SP-04 &
    E4 -- IOTA Network &
    Pihak tidak sah menjalankan node atau validator palsu yang menyamar sebagai partisipan sah jaringan IOTA untuk ikut serta dalam proses konsensus. \\
    \hline

    SP-05 &
    E5 -- IOTA Gas Station &
    Pihak tidak sah menyamar sebagai IOTA Gas Station resmi untuk memproses atau mensponsori transaksi atas nama layanan yang sah. \\
    \hline

    SP-06 &
    P1 -- Hospital Client &
    Aplikasi client palsu memalsukan diri sebagai Hospital Client. \\
    \hline

    SP-07 &
    P2 -- Patient Client &
    Aplikasi client palsu memalsukan diri sebagai Patient Client. \\
    \hline

    SP-08 &
    P3 -- Ministry Client &
    Aplikasi client palsu memalsukan diri sebagai Ministry Client. \\
    \hline

    SP-09 &
    P4 -- Proxy Re-encryption API &
    \textit{api} luar memalsukan diri sebagai Proxy Re-encryption API. \\
    \hline

    SP-10 &
    P5 -- Authentication \& Access Control &
    Sebuah \textit{service} luar meniru proses Authentication \& Access Control. \\
    \hline

    SP-11 &
    P6 -- Medical Record Management &
    Sebuah sistem luar menyamar sebagai Medical Record Management. \\
    \hline

    SP-12 &
    P7 -- Proxy Re-encryption Engine & Sebuah sistem luar menyamar sebagai Proxy Re-encryption Engine. \\ 
    \hline

    SP-13 & P8 -- IOTA Transaction Client & Sebuah sistem luar memalsukan diri sebagai IOTA Transaction Client. \\
    \hline

\end{xltabular}

Hasil analisis menunjukkan bahwa ancaman spoofing pada sistem DecMed terutama berkaitan dengan pemalsuan identitas pengguna atau layanan yang terlibat dalam proses autentikasi, pemberian akses, dan pemrosesan transaksi. Oleh karena itu, audit trail perlu mampu mencatat bukti autentikasi serta identitas entitas yang melakukan setiap operasi sehingga setiap aksi dapat ditelusuri kembali kepada pelaku yang sebenarnya.



% % STRIDE Tampering    
\subsubsection{Tampering}

Berdasarkan pemetaan STRIDE-per-Element, kategori Tampering dianalisis pada elemen Process, Data Store, dan Data Flow. Analisis ini bertujuan untuk mengidentifikasi kemungkinan terjadinya perubahan terhadap logika proses, data yang disimpan, atau data yang dikirimkan selama komunikasi antarelemen.

\begin{xltabular}{\textwidth}{|>{\centering\arraybackslash}p{1.8cm}|>{\raggedright\arraybackslash}p{3.2cm}|X|}
    \caption{Analisis STRIDE Kategori Tampering}
    \label{tab:stride-tampering} \\
    \hline
    \textbf{Kode Ancaman} &
    \textbf{Komponen} &
    \textbf{Deskripsi Ancaman} \\
    \hline
    \endfirsthead

    \multicolumn{3}{c}%
    {} \\
    \hline
    \textbf{Kode Ancaman} &
    \textbf{Komponen} &
    \textbf{Deskripsi Ancaman} \\
    \hline
    \endhead

    \hline
    \multicolumn{3}{r}{} \\
    \endfoot

    \hline
    \endlastfoot

    TP-01 &
    P1 -- Hospital Client &
    Logika atau konfigurasi aplikasi Hospital Client dimodifikasi oleh pihak yang tidak berwenang sehingga permintaan akses atau data yang ditampilkan kepada tenaga medis dapat dimanipulasi. \\
    \hline

    TP-02 &
    P2 -- Patient Client &
    Logika atau konfigurasi aplikasi Patient Client dimodifikasi oleh pihak yang tidak berwenang sehingga parameter persetujuan akses (\textit{consent}) dapat diubah sebelum dikirim ke backend. \\
    \hline

    TP-03 &
    P3 -- Ministry Client &
    Logika atau konfigurasi aplikasi Ministry Client dimodifikasi oleh pihak yang tidak berwenang sehingga data registrasi fasilitas atau tenaga medis dapat dipalsukan sebelum dikirim ke ledger. \\
    \hline

    TP-04 &
    P4 -- Proxy Re-encryption API &
    Logika atau konfigurasi Proxy Re-encryption API dimodifikasi oleh pihak yang tidak berwenang sehingga validasi parameter atau otorisasi permintaan tidak lagi berjalan sebagaimana mestinya. \\
    \hline

    TP-05 &
    P5 -- Authentication \& Access Control &
    Logika atau konfigurasi autentikasi dan otorisasi dimodifikasi oleh pihak yang tidak berwenang sehingga verifikasi \textit{signature}, \textit{nonce}, atau validasi JWT tidak lagi berjalan sesuai mekanisme yang seharusnya. \\
    \hline

    TP-06 &
    P6 -- Medical Record Management &
    Logika atau konfigurasi Medical Record Management dimodifikasi oleh pihak yang tidak berwenang sehingga proses pembuatan atau pembaruan rekam medis tidak lagi mengikuti aturan bisnis yang ditetapkan. \\
    \hline

    TP-07 &
    P7 -- Proxy Re-encryption Engine &
    Logika atau konfigurasi proses re-enkripsi dimodifikasi oleh pihak yang tidak berwenang sehingga menghasilkan \textit{ciphertext} yang dapat diakses oleh pihak yang tidak sesuai dengan permintaan akses yang sah. \\
    \hline

    TP-08 &
    P8 -- IOTA Transaction Client &
    Logika atau konfigurasi pembentukan, penandatanganan, atau pengiriman transaksi dimodifikasi oleh pihak yang tidak berwenang sehingga transaksi yang dikirim ke jaringan IOTA tidak lagi sesuai dengan permintaan pengguna. \\
    \hline

    TP-09 &
    D1 -- Redis Cache &
    Data yang tersimpan pada Redis Cache, seperti \textit{nonce} atau \textit{kfrag}, dimodifikasi oleh pihak yang tidak berwenang sehingga menyebabkan proses autentikasi atau re-enkripsi gagal atau salah sasaran. \\
    \hline

    TP-10 &
    D2 -- IPFS Cluster Storage &
    Data yang tersimpan pada IPFS dimodifikasi oleh pihak yang tidak berwenang sehingga integritas data tidak lagi sesuai dengan CID yang tercatat pada ledger. \\
    \hline

    TP-11 &
    D3 -- Client Key Store &
    \textit{Key} yang tersimpan pada perangkat dimodifikasi oleh pihak yang tidak berwenang sehingga proses dekripsi, penandatanganan, atau re-enkripsi menghasilkan keluaran yang tidak sesuai. \\
    \hline

    TP-12 &
    D4 -- IOTA On-Chain State &
    Meskipun bersifat \textit{immutable}, kerentanan pada logika \textit{smart contract} memungkinkan data yang dikomit ke ledger IOTA tidak sesuai dengan aturan atau kondisi yang seharusnya diterapkan. \\
    \hline

    TP-13 &
    F1 -- Patient PIN \& Session &
    Nilai PIN atau session dimodifikasi \textit{on transit} oleh pihak yang tidak berwenang sehingga proses autentikasi pasien gagal atau salah sasaran. \\
    \hline

    TP-14 &
    F2 -- QR Access Delegation Payload &
    Payload QR (\textit{iota\_address} dan \textit{pre\_public\_key}) dimodifikasi \textit{on transit} oleh pihak yang tidak berwenang sehingga delegasi akses diberikan kepada pihak yang bukan tenaga medis yang sebenarnya. \\
    \hline

    TP-15 &
    F3 -- Hospital Personnel PIN &
    Nilai PIN tenaga medis yang dikirim ke Hospital Client dimodifikasi oleh pihak yang tidak berwenang sehingga proses \textit{sign-in} tidak sesuai dengan identitas asli. \\
    \hline

    TP-16 &
    F4 -- Ministry Admin Credentials &
    Kredensial administrator Kementerian Kesehatan yang dikirim ke Ministry Client dimodifikasi oleh pihak yang tidak berwenang sehingga memungkinkan aksi registrasi dilakukan dengan hak yang tidak sesuai. \\
    \hline

    TP-17 &
    F5 -- IOTA Key Material (ke Patient Client) &
    Materi kunci yang diambil dari Client Key Store dimodifikasi oleh pihak yang tidak berwenang sehingga menyebabkan kegagalan atau kesalahan proses kriptografis pasien. \\
    \hline

    TP-18 &
    F6 -- IOTA Key Material (ke Hospital Client) &
    Materi kunci yang diambil dari Client Key Store dimodifikasi oleh pihak yang tidak berwenang sehingga menyebabkan kegagalan atau kesalahan proses kriptografis tenaga medis. \\
    \hline

    TP-19 &
    F7 -- Proxy API Request (JWT) &
    Token JWT pada permintaan dari Hospital Client menuju Proxy Re-encryption API dimodifikasi oleh pihak yang tidak berwenang sehingga klaim \textit{role} atau \textit{access type} dapat diubah. \\
    \hline

    TP-20 &
    F8 -- Access Delegation API Call &
    Parameter pada panggilan delegasi akses dari Patient Client dimodifikasi oleh pihak yang tidak berwenang sehingga durasi atau cakupan akses tidak sesuai dengan yang disetujui pasien. \\
    \hline

    TP-21 &
    F9 -- On-Chain Access Grant Transaction &
    Payload transaksi pemberian hak akses dimodifikasi oleh pihak yang tidak berwenang sehingga durasi atau cakupan akses tidak sesuai dengan persetujuan pasien sebelum dikomit ke ledger. \\
    \hline

    TP-22 &
    F10 -- Hospital Registration Transaction &
    Data registrasi fasilitas kesehatan atau tenaga medis dimodifikasi oleh pihak yang tidak berwenang sehingga entitas yang terdaftar tidak sesuai dengan permintaan administrator yang asli. \\
    \hline

    TP-23 &
    F11 -- API Route Dispatch (ke P5) &
    Instruksi \textit{routing} dari Proxy Re-encryption API menuju Authentication \& Access Control dimodifikasi oleh pihak yang tidak berwenang sehingga permintaan diarahkan ke logika validasi yang salah. \\
    \hline

    TP-24 &
    F12 -- API Route Dispatch (ke P6) &
    Instruksi \textit{routing} dari Proxy Re-encryption API menuju Medical Record Management dimodifikasi oleh pihak yang tidak berwenang sehingga permintaan diarahkan ke logika pemrosesan yang salah. \\
    \hline

    TP-25 &
    F13 -- Authentication Nonce &
    Nilai \textit{nonce} yang dipertukarkan antara Authentication \& Access Control dan IOTA Transaction Client dimodifikasi oleh pihak yang tidak berwenang sehingga memungkinkan serangan \textit{replay} atau \textit{bypass} validasi. \\
    \hline

    TP-26 &
    F14 -- Proxy Access Keys (ke P7) &
    Kunci akses proxy yang diteruskan menuju Proxy Re-encryption Engine dimodifikasi oleh pihak yang tidak berwenang sehingga proses re-enkripsi menghasilkan \textit{ciphertext} yang salah sasaran. \\
    \hline

    TP-27 &
    F15 -- JWT Access Token &
    Token akses yang dikembalikan kepada client dimodifikasi oleh pihak yang tidak berwenang sehingga client menggunakan klaim akses yang tidak sah. \\
    \hline

    TP-28 &
    F16 -- Re-encryption Job &
    Parameter atau instruksi proses re-enkripsi yang dikirim dari Medical Record Management ke Proxy Re-encryption Engine dimodifikasi oleh pihak yang tidak berwenang sehingga hasil re-enkripsi tidak sesuai dengan permintaan akses yang sah. \\
    \hline

    TP-29 &
    F17 -- Proxy Access Keys (ke P8) &
    Kunci akses hasil proses re-enkripsi yang dikembalikan ke IOTA Transaction Client dimodifikasi oleh pihak yang tidak berwenang sehingga memengaruhi keputusan transaksi selanjutnya. \\
    \hline

    TP-30 &
    F18 -- Encrypted Medical Blob &
    Blob data medis terenkripsi dimodifikasi oleh pihak yang tidak berwenang sehingga hasil re-enkripsi tidak sesuai dengan data klinis yang sebenarnya. \\
    \hline

    TP-31 &
    F19 -- Encrypted Medical Blob (ke Storage) &
    Blob data klinis terenkripsi yang dikirim menuju penyimpanan dimodifikasi oleh pihak yang tidak berwenang sehingga data yang tersimpan tidak sesuai dengan yang dienkripsi pada awalnya. \\
    \hline

    TP-32 &
    F20 -- Ciphertext Fragment &
    Fragmen \textit{ciphertext} yang dikirim ke client untuk didekripsi secara lokal dimodifikasi oleh pihak yang tidak berwenang sehingga proses dekripsi gagal atau menghasilkan data yang salah. \\
    \hline

    TP-33 &
    F21 -- On-Chain Metadata (P8 ke D4) &
    Metadata yang dikirim oleh IOTA Transaction Client untuk disimpan \textit{on-chain} dimodifikasi oleh pihak yang tidak berwenang sehingga metadata pada ledger tidak sesuai dengan transaksi asli. \\
    \hline

    TP-34 &
    F22 -- On-Chain Metadata (D4 ke P6) &
    Metadata yang diambil dari ledger untuk diproses oleh Medical Record Management dimodifikasi oleh pihak yang tidak berwenang sehingga keputusan proses didasarkan pada data yang salah. \\
    \hline

    TP-35 &
    F23 -- Move Transaction &
    Transaksi Move dimodifikasi oleh pihak yang tidak berwenang sebelum divalidasi oleh \textit{smart contract} sehingga transaksi yang diproses tidak sesuai dengan yang dikirim oleh sistem. \\
    \hline

    TP-36 &
    F24 -- Gas Sponsorship Request &
    Permintaan sponsor \textit{gas} dari IOTA Transaction Client menuju Gas Station dimodifikasi oleh pihak yang tidak berwenang sehingga sponsor mendanai transaksi yang bukan sebenarnya diminta. \\
    \hline

    TP-37 &
    F25 -- Gas Coin Reservation &
    Konfirmasi reservasi \textit{gas coin} dari Gas Station dimodifikasi oleh pihak yang tidak berwenang sehingga menyebabkan transaksi gagal atau salah alokasi. \\
    \hline

    TP-38 &
    F26 -- Personnel Role Lookup &
    Hasil pencarian \textit{role} tenaga medis yang digunakan oleh Authentication \& Access Control dimodifikasi oleh pihak yang tidak berwenang sehingga keputusan otorisasi berbasis \textit{role} menjadi salah. \\
    \hline

    TP-39 &
    F27 -- Medical Record Metadata &
    Metadata rekam medis yang dipertukarkan antara proses backend dan penyimpanan dimodifikasi oleh pihak yang tidak berwenang sehingga referensi atau lokasi data klinis tidak sesuai dengan yang sebenarnya. \\
    \hline

    TP-40 &
    F28 -- Client-Side Decryption &
    Instruksi atau parameter dekripsi lokal pada sisi client dimodifikasi oleh pihak yang tidak berwenang, misalnya melalui aplikasi yang di-\textit{patch}, sehingga data yang ditampilkan kepada tenaga medis tidak sesuai dengan data asli. \\
    \hline

    TP-41 &
    F29 -- Access Log View &
    Permintaan atau hasil tampilan log akses dimodifikasi oleh pihak yang tidak berwenang sehingga log yang dilihat tidak mencerminkan riwayat akses yang sebenarnya terjadi. \\
    \hline

\end{xltabular}

Hasil analisis menunjukkan bahwa ancaman tampering tidak hanya terjadi pada data yang dikirimkan melalui data yang sudah tersimpan di data store, tetapi juga pada modifikasi logika proses serta data yang sedang dalam proses. Hal ini menunjukkan bahwa audit trail perlu menyediakan bukti yang dapat digunakan untuk memverifikasi integritas data atau aktivitas yang memengaruhi perubahan data selama siklus pemrosesan.

% % STRIDE Repudiation
\subsubsection{Repudiation}

Berdasarkan pemetaan STRIDE-per-Element, kategori Repudiation dianalisis pada elemen External Entity, Process, dan Data Store yang berperan sebagai media pencatatan. Analisis ini bertujuan untuk mengidentifikasi kondisi ketika suatu pihak dapat menyangkal telah melakukan suatu tindakan karena tidak tersedia bukti yang dapat dipertanggungjawabkan.

\begin{xltabular}{\textwidth}{|>{\centering\arraybackslash}p{1.8cm}|>{\raggedright\arraybackslash}p{3.2cm}|X|}
    \caption{Analisis STRIDE Kategori Repudiation}
    \label{tab:stride-repudiation} \\
    \hline
    \textbf{Kode Ancaman} &
    \textbf{Komponen} &
    \textbf{Deskripsi Ancaman} \\
    \hline
    \endfirsthead

    \multicolumn{3}{c}%
    {} \\
    \hline
    \textbf{Kode Ancaman} &
    \textbf{Komponen} &
    \textbf{Deskripsi Ancaman} \\
    \hline
    \endhead

    \hline
    \multicolumn{3}{r}{} \\
    \endfoot

    \hline
    \endlastfoot

    RP-01 &
    E1 -- Patient &
    Pasien menyangkal pernah memberikan atau mencabut akses kepada tenaga medis tertentu, terutama setelah terjadi insiden atau pada kondisi darurat. \\
    \hline

    RP-02 &
    E2 -- Hospital Personnel &
    Tenaga medis menyangkal pernah mengakses atau memodifikasi rekam medis pasien tertentu, meskipun tindakan tersebut benar-benar dilakukan selama periode akses yang sah. \\
    \hline

    RP-03 &
    E3 -- Ministry Administrator &
    Administrator Kementerian Kesehatan menyangkal pernah me-\textit{register} fasilitas kesehatan tertentu, terutama apabila registrasi tersebut kemudian disalahgunakan. \\
    \hline

    RP-04 &
    E4 -- IOTA Network &
    Validator atau node jaringan tidak menghasilkan bukti atas keterlibatannya dalam memfinalisasi transaksi tertentu sehingga sulit menentukan pihak yang bertanggung jawab atas hasil konsensus. \\
    \hline

    RP-05 &
    E5 -- IOTA Gas Station &
    IOTA Gas Station tidak menghasilkan bukti pernah mensponsori transaksi tertentu, meskipun transaksi tersebut benar-benar menggunakan layanannya. \\
    \hline

    RP-06 &
    P1 -- Hospital Client &
    Hospital Client tidak menghasilkan bukti yang mengikat aktivitas dengan sesi yang sah sehingga tenaga medis dapat menyangkal permintaan yang sebenarnya berasal dari perangkat atau sesinya. \\
    \hline

    RP-07 &
    P2 -- Patient Client &
    Patient Client tidak menghasilkan bukti yang mengikat aktivitas delegasi atau pencabutan akses dengan sesi yang sah sehingga pasien dapat menyangkal tindakan tersebut. \\
    \hline

    RP-08 &
    P3 -- Ministry Client &
    Ministry Client tidak menghasilkan bukti yang mengikat aktivitas registrasi dengan sesi yang sah sehingga administrator dapat menyangkal telah melakukan registrasi tersebut. \\
    \hline

    RP-09 &
    P4 -- Proxy Re-encryption API &
    Tidak terdapat jejak yang menghubungkan setiap permintaan API dengan identitas pemohon yang sebenarnya sehingga sulit membuktikan pihak yang memicu suatu operasi. \\
    \hline

    RP-10 &
    P5 -- Authentication \& Access Control &
    Tidak ada bukti permintaan dan proses pemberian otoritas yang dapat ditelusuri kembali ke konteks tertentu sehingga dasar pemberian atau penolakan akses sulit dibuktikan ketika terjadi sengketa. \\
    \hline

    RP-11 &
    P6 -- Medical Record Management &
    Tenaga medis dapat menyangkal telah membuat atau mengubah entri rekam medis tertentu apabila tidak terdapat bukti yang mengikat tindakan tersebut dengan identitas pelaku dan waktu pelaksanaannya. \\
    \hline

    RP-12 &
    P7 -- Proxy Re-encryption Engine &
    Tidak terdapat bukti bahwa permintaan dan proses re-enkripsi dilakukan berdasarkan permintaan akses yang sah sehingga legitimasi hasilnya sulit dibuktikan apabila dipersengketakan. \\
    \hline

    RP-13 &
    P8 -- IOTA Transaction Client &
    Tidak terdapat bukti yang mengikat transaksi dengan pihak yang pertama kali memicunya sehingga pengguna dapat menyangkal telah memicu transaksi tersebut. \\
    \hline

    RP-14 &
    D4 -- IOTA On-Chain State &
    Tidak ada pengikatan antara identitas pemohon dan entri transaksi pada ledger sehingga pihak terkait dapat menyangkal keterkaitannya dengan transaksi tersebut, meskipun ledger bersifat \textit{immutable}. \\
    \hline

\end{xltabular}

Hasil analisis menunjukkan bahwa sebagian besar ancaman repudiation muncul akibat tidak adanya bukti yang mengikat antara identitas pelaku, aktivitas yang dilakukan, dan waktu pelaksanaan aktivitas tersebut. Oleh karena itu, audit trail perlu mampu menyediakan bukti yang tidak dapat disangkal (non-repudiation) untuk mendukung proses investigasi atau penyelesaian sengketa.

% % STRIDE Information Disclosure
\subsubsection{Information Disclosure}

Berdasarkan pemetaan STRIDE-per-Element, kategori Information Disclosure dianalisis pada elemen Process, Data Store, dan Data Flow. Analisis ini bertujuan untuk mengidentifikasi kemungkinan terbukanya informasi sensitif selama proses pemrosesan, penyimpanan, atau pertukaran data antar komponen sistem.

\begin{xltabular}{\textwidth}{|>{\centering\arraybackslash}p{1.8cm}|>{\raggedright\arraybackslash}p{3.2cm}|X|}
    \caption{Analisis STRIDE Kategori Information Disclosure}
    \label{tab:stride-information-disclosure} \\
    \hline
    \textbf{Kode Ancaman} &
    \textbf{Komponen} &
    \textbf{Deskripsi Ancaman} \\
    \hline
    \endfirsthead

    \multicolumn{3}{c}%
    {} \\
    \hline
    \textbf{Kode Ancaman} &
    \textbf{Komponen} &
    \textbf{Deskripsi Ancaman} \\
    \hline
    \endhead

    \hline
    \multicolumn{3}{r}{} \\
    \endfoot

    \hline
    \endlastfoot

    ID-01 &
    P1 -- Hospital Client &
    Pihak tidak berwenang melihat data rekam medis yang telah didekripsi dan ditampilkan pada Hospital Client, mengungkap informasi medis pasien. \\
    \hline

    ID-02 &
    P2 -- Patient Client &
    Pihak tidak berwenang mengakses data pribadi pasien melalui Patient Client, seperti identitas, \textit{recovery phrase}, atau riwayat akses. \\
    \hline

    ID-03 &
    P3 -- Ministry Client &
    Pihak tidak berwenang mengakses informasi administrator, fasilitas kesehatan, atau tenaga medis yang dikelola melalui Ministry Client. \\
    \hline

    ID-04 &
    P4 -- Proxy Re-encryption API &
    Informasi sensitif yang diproses oleh Proxy Re-encryption API, seperti token autentikasi atau parameter akses, terekspos kepada pihak yang tidak berwenang. \\
    \hline

    ID-05 &
    P5 -- Authentication \& Access Control &
    Informasi autentikasi dan otorisasi, seperti klaim JWT, \textit{role}, atau hak akses pengguna, terekspos kepada pihak yang tidak berwenang. \\
    \hline

    ID-06 &
    P6 -- Medical Record Management &
    Data rekam medis pasien yang sedang diproses terekspos sebelum dienkripsi atau setelah didekripsi untuk kebutuhan pemrosesan. \\
    \hline

    ID-07 &
    P7 -- Proxy Re-encryption Engine &
    Material kriptografi yang digunakan selama proses re-enkripsi, seperti \textit{kfrag} dan parameter re-enkripsi, terekspos sehingga dapat disalahgunakan untuk mengakses data pasien. \\
    \hline

    ID-08 &
    P8 -- IOTA Transaction Client &
    Metadata transaksi dan parameter akses yang diproses sebelum dikirim ke jaringan IOTA terekspos kepada pihak yang tidak berwenang. \\
    \hline

    ID-09 &
    D1 -- Redis Cache &
    Data sementara yang tersimpan pada Redis Cache, seperti \textit{nonce}, \textit{kfrag}, atau token sesi, diakses oleh pihak yang tidak berwenang. \\
    \hline

    ID-10 &
    D2 -- IPFS Cluster Storage &
    Blob data medis yang tersimpan pada IPFS dapat diakses oleh siapa pun yang mengetahui CID sehingga kerahasiaannya sepenuhnya bergantung pada kekuatan mekanisme enkripsi yang digunakan. \\
    \hline

    ID-11 &
    D3 -- Client Key Store &
    \textit{Key} yang tersimpan pada Client Key Store diakses oleh pihak yang tidak berwenang sehingga membahayakan kerahasiaan data pasien. \\
    \hline

    ID-12 &
    D4 -- IOTA On-Chain State &
    Metadata transaksi yang tersimpan pada ledger publik mengungkap pola akses, hubungan antara pasien dan tenaga medis, atau aktivitas sistem melalui analisis metadata. \\
    \hline

    ID-13 &
    F1 -- Patient PIN \& Session &
    PIN atau token sesi pasien terekspos selama transmisi sehingga dapat digunakan oleh pihak yang tidak berwenang. \\
    \hline

    ID-14 &
    F2 -- QR Access Delegation Payload &
    Payload QR yang berisi \textit{iota\_address} dan \textit{pre\_public\_key} tenaga medis terekspos kepada pihak selain pasien yang dituju. \\
    \hline

    ID-15 &
    F3 -- Hospital Personnel PIN &
    PIN tenaga medis terekspos selama transmisi menuju Hospital Client. \\
    \hline

    ID-16 &
    F4 -- Ministry Admin Credentials &
    Kredensial administrator Kementerian Kesehatan terekspos selama transmisi menuju Ministry Client. \\
    \hline

    ID-17 &
    F5 -- IOTA Key Material (Patient Client) &
    Material kunci yang dipindahkan dari Client Key Store ke Patient Client terekspos selama proses penggunaan. \\
    \hline

    ID-18 &
    F6 -- IOTA Key Material (Hospital Client) &
    Material kunci yang dipindahkan dari Client Key Store ke Hospital Client terekspos selama proses penggunaan. \\
    \hline

    ID-19 &
    F7 -- Proxy API Request (JWT) &
    Token JWT yang dikirim menuju Proxy Re-encryption API terekspos selama transmisi. \\
    \hline

    ID-20 &
    F8 -- Access Delegation API Call &
    Parameter delegasi akses, termasuk cakupan dan durasi akses, terekspos selama transmisi menuju backend. \\
    \hline

    ID-21 &
    F9 -- On-Chain Access Grant Transaction &
    Metadata transaksi pemberian hak akses mengekspos hubungan antara pasien dan tenaga medis ketika di-\textit{commit} ke ledger. \\
    \hline

    ID-22 &
    F10 -- Hospital Registration Transaction &
    Informasi registrasi fasilitas kesehatan atau tenaga medis terekspos selama transmisi menuju ledger. \\
    \hline

    ID-23 &
    F11 -- API Route Dispatch (ke P5) &
    Informasi \textit{routing} internal mengungkap detail permintaan autentikasi kepada pihak yang tidak berwenang. \\
    \hline

    ID-24 &
    F12 -- API Route Dispatch (ke P6) &
    Informasi \textit{routing} internal mengungkap detail permintaan rekam medis kepada pihak yang tidak berwenang. \\
    \hline

    ID-25 &
    F13 -- Authentication Nonce &
    Nilai \textit{nonce} autentikasi terekspos selama transmisi dan dimanfaatkan untuk analisis pola autentikasi. \\
    \hline

    ID-26 &
    F14 -- Proxy Access Keys (ke P7) &
    Kunci akses proxy yang dikirim menuju Proxy Re-encryption Engine terekspos kepada pihak yang tidak berwenang. \\
    \hline

    ID-27 &
    F15 -- JWT Access Token &
    Token akses yang dikembalikan kepada client terekspos selama transmisi. \\
    \hline

    ID-28 &
    F16 -- Re-encryption Job &
    Parameter proses re-enkripsi mengungkap metadata atau informasi akses yang sedang diproses. \\
    \hline

    ID-29 &
    F17 -- Proxy Access Keys (ke P8) &
    Kunci akses hasil proses re-enkripsi yang dikirim menuju IOTA Transaction Client terekspos kepada pihak yang tidak berwenang. \\
    \hline

    ID-30 &
    F18 -- Encrypted Medical Blob &
    Blob data medis terenkripsi yang dikirim selama proses re-enkripsi mengungkap metadata atau pola akses meskipun isi datanya tetap terenkripsi. \\
    \hline

    ID-31 &
    F19 -- Encrypted Medical Blob (ke Storage) &
    Blob data medis terenkripsi yang dikirim menuju media penyimpanan mengungkap metadata atau pola penggunaan data. \\
    \hline

    ID-32 &
    F20 -- Ciphertext Fragment &
    Fragmen \textit{ciphertext} yang dikirim menuju client terekspos sebelum proses dekripsi dilakukan. \\
    \hline

    ID-33 &
    F21 -- On-Chain Metadata (P8 ke D4) &
    Metadata yang dikirim ke ledger mengungkap pola transaksi dan hubungan antar entitas karena sifat ledger yang transparan. \\
    \hline

    ID-34 &
    F22 -- On-Chain Metadata (D4 ke P6) &
    Metadata yang diambil dari ledger terekspos selama transmisi menuju Medical Record Management. \\
    \hline

    ID-35 &
    F23 -- Move Transaction &
    Informasi transaksi Move terekspos selama transmisi menuju jaringan IOTA. \\
    \hline

    ID-36 &
    F24 -- Gas Sponsorship Request &
    Permintaan sponsor \textit{gas} mengungkap metadata transaksi atau pola penggunaan layanan kepada IOTA Gas Station. \\
    \hline

    ID-37 &
    F25 -- Gas Coin Reservation &
    Informasi reservasi \textit{gas coin} mengungkap alokasi sumber daya atau aktivitas transaksi. \\
    \hline

    ID-38 &
    F26 -- Personnel Role Lookup &
    Informasi \textit{role} tenaga medis mengungkap struktur hak akses dalam suatu institusi. \\
    \hline

    ID-39 &
    F27 -- Medical Record Metadata &
    Metadata rekam medis mengungkap referensi atau lokasi penyimpanan data pasien tanpa membuka isi rekam medis. \\
    \hline

    ID-40 &
    F28 -- Client-Side Decryption &
    Data rekam medis hasil dekripsi pada sisi client terekspos sebelum digunakan oleh tenaga medis. \\
    \hline

    ID-41 &
    F29 -- Access Log View &
    Informasi riwayat akses mengungkap siapa mengakses data siapa, waktu akses, atau pola penggunaan sistem kepada pihak yang tidak berwenang. \\
    \hline

\end{xltabular}

Hasil analisis menunjukkan bahwa ancaman kebocoran informasi tidak hanya berasal dari isi data klinis, tetapi juga dari metadata transaksi, parameter akses, atau material kriptografi yang digunakan sistem. Hal ini menunjukkan bahwa audit trail perlu dirancang agar tetap menyediakan informasi yang diperlukan untuk investigasi tanpa mengungkap data sensitif yang tidak diperlukan.

% % STRIDE Denial of Service
\subsubsection{Denial of Service}

Berdasarkan pemetaan STRIDE-per-Element, kategori Denial of Service dianalisis pada elemen Process, Data Store, dan Data Flow. Namun analisis pada penelitian ini dibatasi pada elemen Process dan External Entity, tanpa mencakup elemen Data Flow. Hal ini didasarkan pada pertimbangan bahwa ancaman DoS pada Data Flow pada dasarnya merupakan manifestasi dari pembebanan berlebih terhadap Process atau Data Store pada titik akhir aliran data tersebut, sehingga telah tercakup secara implisit dalam analisis pada elemen Process. Selain itu, karena penelitian ini berfokus pada kebutuhan audit trail berbasis pencatatan peristiwa (\textit{event-based logging}), sementara mitigasi DoS umumnya memerlukan mekanisme pemantauan volume dan laju permintaan (\textit{rate limiting}, \textit{traffic monitoring}) yang berada di luar cakupan audit trail, maka analisis DoS pada level Data Flow tidak menghasilkan kebutuhan audit trail yang berbeda secara signifikan dari analisis pada elemen Process.

\begin{xltabular}{\textwidth}{|>{\centering\arraybackslash}p{1.8cm}|>{\raggedright\arraybackslash}p{3.2cm}|X|}
    \caption{Analisis STRIDE Kategori Denial of Service}
    \label{tab:stride-denial-of-service} \\
    \hline
    \textbf{Kode Ancaman} &
    \textbf{Komponen} &
    \textbf{Deskripsi Ancaman} \\
    \hline
    \endfirsthead

    \multicolumn{3}{c}%
    {} \\
    \hline
    \textbf{Kode Ancaman} &
    \textbf{Komponen} &
    \textbf{Deskripsi Ancaman} \\
    \hline
    \endhead

    \hline
    \multicolumn{3}{r}{} \\
    \endfoot

    \hline
    \endlastfoot

    DS-01 &
    P1 -- Hospital Client &
    Permintaan API dalam jumlah besar dikirimkan secara berulang menuju Hospital Client sehingga membebani kemampuan aplikasi memproses permintaan tenaga medis yang sah. \\
    \hline

    DS-02 &
    P2 -- Patient Client &
    Permintaan delegasi atau pencabutan akses dikirimkan secara berulang dalam volume tinggi menuju Patient Client sehingga membebani kemampuan aplikasi memproses permintaan pasien yang sah. \\
    \hline

    DS-03 &
    P3 -- Ministry Client &
    Permintaan registrasi fasilitas kesehatan atau tenaga medis dikirimkan secara berulang dalam volume tinggi menuju Ministry Client sehingga menghambat pemrosesan registrasi yang sah. \\
    \hline

    DS-04 &
    P4 -- Proxy Re-encryption API &
    Permintaan yang tidak sah atau berulang secara masif dikirimkan ke Proxy Re-encryption API. \\
    \hline

    DS-05 &
    P5 -- Authentication \& Access Control &
    Permintaan autentikasi atau validasi token dikirimkan secara berulang dalam volume tinggi. \\
    \hline

    DS-06 &
    P6 -- Medical Record Management &
    Permintaan pembuatan atau pembaruan rekam medis dikirimkan secara berulang dalam volume tinggi. \\
    \hline

    DS-07 &
    P7 -- Proxy Re-encryption Engine &
    Permintaan proses re-enkripsi dikirimkan secara berulang atau dengan parameter yang menyebabkan komputasi kriptografi menjadi sangat berat. \\
    \hline

    DS-08 &
    P8 -- IOTA Transaction Client &
    Permintaan pembentukan atau pengiriman transaksi dikirimkan secara berulang dalam volume tinggi menuju IOTA Transaction Client sehingga membebani kemampuannya memproses transaksi yang sah. \\
    \hline

    DS-09 &
    D1 -- Redis Cache &
    Penulisan data sementara dilakukan secara berulang dalam volume tinggi. \\
    \hline

    DS-10 &
    D2 -- IPFS Cluster Storage &
    Unggahan blob data medis dilakukan secara berulang dalam volume tinggi secara sengaja. \\
    \hline

    DS-11 &
    E5 -- IOTA Gas Station &
    Permintaan sponsor gas dikirimkan secara berulang dalam volume tinggi. \\
    \hline

\end{xltabular}

Hasil analisis menunjukkan bahwa ancaman denial of service pada sistem DecMed umumnya terjadi akibat pembebanan berlebih pada proses atau penyimpanan data sementara melalui pengiriman permintaan secara masif. Meskipun kategori ini tidak dianalisis secara mendalam pada penelitian ini karena berada di luar fokus perancangan kebutuhan audit trail, temuan awal ini dapat menjadi masukan bagi penelitian atau pengembangan selanjutnya terkait mekanisme pembatasan laju permintaan (\textit{rate limiting}) dan pemantauan ketersediaan layanan.

% % STRIDE Elevation of Privilege
\subsubsection{Elevation of Privilege}

Berdasarkan pemetaan STRIDE-per-Element, kategori Elevation of Privilege dianalisis pada elemen Process. Analisis ini bertujuan untuk mengidentifikasi kemungkinan pengguna atau proses memperoleh hak akses yang melebihi kewenangan yang diberikan akibat kelemahan pada mekanisme autentikasi, otorisasi, atau logika bisnis sistem.

\begin{xltabular}{\textwidth}{|>{\centering\arraybackslash}p{1.8cm}|>{\raggedright\arraybackslash}p{3.2cm}|X|}
    \caption{Analisis STRIDE Kategori Elevation of Privilege}
    \label{tab:stride-elevation-of-privilege} \\
    \hline
    \textbf{Kode Ancaman} &
    \textbf{Komponen} &
    \textbf{Deskripsi Ancaman} \\
    \hline
    \endfirsthead

    \multicolumn{3}{c}%
    {} \\
    \hline
    \textbf{Kode Ancaman} &
    \textbf{Komponen} &
    \textbf{Deskripsi Ancaman} \\
    \hline
    \endhead

    \hline
    \multicolumn{3}{r}{} \\
    \endfoot

    \hline
    \endlastfoot

    EP-01 &
    P1 -- Hospital Client &
    Kerentanan pada Hospital Client dieksploitasi sehingga pengguna memperoleh akses terhadap fungsi atau fitur yang seharusnya hanya tersedia bagi pengguna dengan hak akses yang lebih tinggi, seperti fungsi administratif. \\
    \hline

    EP-02 &
    P2 -- Patient Client &
    Kerentanan pada Patient Client dieksploitasi sehingga pasien memperoleh akses terhadap sumber daya atau fungsi yang berada di luar kewenangannya, seperti mengakses data pasien lain atau mengubah konfigurasi sistem. \\
    \hline

    EP-03 &
    P3 -- Ministry Client &
    Kerentanan pada Ministry Client dieksploitasi sehingga administrator memperoleh kemampuan yang melebihi kewenangannya, seperti mengubah kebijakan atau konfigurasi keamanan sistem tanpa otorisasi yang semestinya. \\
    \hline

    EP-04 &
    P4 -- Proxy Re-encryption API &
    Kerentanan pada Proxy Re-encryption API dieksploitasi sehingga pengguna dapat melewati mekanisme otorisasi dan mengakses proses backend secara langsung tanpa melalui validasi Authentication \& Access Control. \\
    \hline

    EP-05 &
    P5 -- Authentication \& Access Control &
    Kesalahan pada mekanisme autentikasi atau otorisasi memungkinkan pengguna memperoleh hak akses yang lebih tinggi daripada yang diberikan, misalnya akibat kelemahan validasi role, klaim JWT, atau \textit{signature}. \\
    \hline

    EP-06 &
    P6 -- Medical Record Management &
    Kerentanan pada Medical Record Management dieksploitasi sehingga pengguna dapat mengakses atau mengubah rekam medis di luar cakupan capability atau hak akses yang dimilikinya. \\
    \hline

    EP-07 &
    P7 -- Proxy Re-encryption Engine &
    Kerentanan pada Proxy Re-encryption Engine dieksploitasi sehingga proses re-enkripsi dapat dilakukan untuk penerima yang tidak berwenang, memberikan kemampuan mengakses data kepada pihak yang seharusnya tidak memiliki hak tersebut. \\
    \hline

    EP-08 &
    P8 -- IOTA Transaction Client &
    Kerentanan pada IOTA Transaction Client dieksploitasi sehingga transaksi dapat dibentuk, ditandatangani, atau dikirim menggunakan capability atau hak akses yang melebihi otorisasi yang diberikan, atau melewati mekanisme validasi CapBAC pada smart contract. \\
    \hline

\end{xltabular}

Hasil analisis menunjukkan bahwa sebagian besar ancaman elevation of privilege berkaitan dengan kegagalan penegakan mekanisme Capability-Based Access Control (CapBAC) pada berbagai proses di dalam sistem. Oleh karena itu, audit trail perlu mampu mencatat informasi mengenai hak akses yang diberikan, hak akses yang digunakan, serta hasil keputusan otorisasi agar setiap penyalahgunaan privilege dapat diidentifikasi dan ditelusuri.

Berdasarkan hasil identifikasi ancaman menggunakan STRIDE, dapat diamati bahwa beberapa pola ancaman muncul secara berulang pada berbagai elemen sistem. Misalnya, ancaman spoofing menekankan pentingnya pencatatan identitas dan hasil autentikasi, sedangkan ancaman repudiation menekankan perlunya bukti yang mengikat antara identitas, aktivitas, dan waktu kejadian. Demikian pula, ancaman tampering, information disclosure, dan elevation of privilege menunjukkan kebutuhan untuk mencatat informasi mengenai integritas data, akses terhadap informasi sensitif, serta penggunaan hak akses selama proses berlangsung. Pola-pola ancaman tersebut kemudian digunakan sebagai dasar dalam menyusun kebutuhan konten audit trail yang dibahas pada subbab berikutnya.